%%%%%%%%%%%%%%%%%%%%%%%%%%%%%%%%%%
%    DO NOT MODIFY THIS CODE     %
%%%%%%%%%%%%%%%%%%%%%%%%%%%%%%%%%%
\documentclass[twoside, 11pt]{article}


\usepackage[english]{babel}
\usepackage{csquotes}


\usepackage{geometry}
\geometry{a4paper,left=2.5cm,right=2cm,top=2cm,bottom=2.5cm, includehead, headheight=65pt}

\usepackage{tinos}
\usepackage[T1]{fontenc}
\usepackage{josefin}
\usepackage{setspace}
\newcommand{\josefinRegular}[1]{{\fontfamily{JosefinSans-LF}\selectfont{#1}}}
\newcommand{\josefinSemiBold}[1]{\josefinsemibold{#1}}
\newcommand{\josefinBold}[1]{\textbf{\josefinRegular{#1}}}
\newcommand{\josefinItalic}[1]{\textit{\josefinRegular{#1}}}

\usepackage[style=apa]{biblatex}
\usepackage{xcolor}
\usepackage[hidelinks]{hyperref}

\makeatletter
\newcommand{\ArticleTitlePTBR}[1]{\def\@ArticleTitlePTBR{#1}}
\newcommand{\ArticleTitleENG}[1]{\def\@ArticleTitleENG{#1}}
\newcommand{\AbstractPTBR}[1]{\def\@AbstractPTBR{#1}}
\newcommand{\AbstractENG}[1]{\def\@AbstractENG{#1}}
\newcommand{\KeywordsPTBR}[1]{\def\@KeywordsPTBR{#1}}
\newcommand{\KeywordsENG}[1]{\def\@KeywordsENG{#1}}



%%%%%%%%%%%%%%%%%%%%%%%%%%%%%%%%%%
%    DO NOT MODIFY THIS CODE     %
%%%%%%%%%%%%%%%%%%%%%%%%%%%%%%%%%%

\renewcommand{\maketitle}{
    \noindent
    \setstretch{1.5}
    \begin{flushleft}
        \medskip
        \noindent{\fontsize{19}{0}{\josefinBold{\@ArticleTitleENG}}}
    
        \bigskip

        \noindent{\fontsize{17}{0}{\josefinBold{\@ArticleTitlePTBR}}}

        \noindent\fontsize{10}{0}{\@author}
    \end{flushleft}
    \setstretch{1.0}
    \vspace{-0.4cm}
  \@thanks
}

%%%%%%%%%%%%%%%%%%%%%%%%%%%%%%%%%%
%    DO NOT MODIFY THIS CODE     %
%%%%%%%%%%%%%%%%%%%%%%%%%%%%%%%%%%

\newcommand{\makeAbstracts}
{
    \noindent\fontsize{11}{0}{{\josefinBold{ABSTRACT}}}
    
    \medskip
    \noindent\small{\@AbstractENG}
    \medskip

    \noindent{\josefinBold{keywords}} \quad \small{\@KeywordsENG}
    
    \bigskip\medskip
    
    \noindent\fontsize{11}{0}{{\josefinBold{RESUMO}}}
    
    \medskip
    \noindent\small{\@AbstractPTBR}
    \medskip
    
    \noindent{\josefinBold{palavras-chave}} \quad \small{\@KeywordsPTBR}
}
\makeatother

%%%%%%%%%%%%%%%%%%%%%%%%%%%%%%%%%%
%    DO NOT MODIFY THIS CODE     %
%%%%%%%%%%%%%%%%%%%%%%%%%%%%%%%%%%

\newcommand{\Init}{%
    \onecolumn
    \setcounter{page}{1}
    \maketitle
    \noindent {\scriptsize \textit{Received: 00 Month 0000 \hfill Received in revised form: 00 Month 0000 \hfill Accepted: 00 Month 0000 \hfill Available online: 00 Month 0000}}

    \bigskip
    
    \makeAbstracts
    \vfill
    
    \clearpage
    \restoregeometry
    
    \newgeometry{left=2.5cm,right=2cm,top=2cm,bottom=2.5cm}
    \linespread{1.5}
}


%%%%%%%%%%%%%%%%%%%%%%%%%%%%%%%%%%
%    DO NOT MODIFY THIS CODE     %
%%%%%%%%%%%%%%%%%%%%%%%%%%%%%%%%%%































%%%%% BIBLIOGRAPHY FILE %%%%%%%%%%%%%%%%%%%%%%%%%%%%%%
\bibliography{references}


%%%%% CUSTOM PACKAGE %%%%%%%%%%%%%%%%%%%%%%%%%%%%%%%%%
\usepackage{amsmath}



%%%%% AUTHORS %%%%%%%%%%%%%%%%%%%%%%%%%%%%%%%%%%%%%%%%
\author{\footnotesize{Author1 (name) -  Prof. Dr., Dpto of Mathematics, UEL, Londrina, Pr, Brazil. E-mail: docente@uel.br (shortly inform their biography)}; \\
%
\footnotesize{Author2 (name) -  Prof. Dr., Dpto of Mathematics, UEL, Londrina, Pr, Brazil. E-mail: docente@uel.br (shortly inform their biography)}\\
}
%%%%%%%%%%%%%%%%%%%%%%%%%%%%%%%%%%%%%%%%%%%%%%%%%%%%%%

\ArticleTitleENG{The title must be concise (not containing more than 100 characters), informative, and must not include a subtitle}
\ArticleTitlePTBR{O título deve ser conciso (não mais que 100 caracteres), informativo e não deve incluir legenda}
\AbstractENG{%
    The abstract must be informative and not only indicate the general scope of the article but also the main results obtained, the utilized methods, the value of the work, and the conclusions drawn. Table numbers, figure numbers, references, or mathematical expressions should not be included. Not exceed 200 words. Once the abstract is finished, it is necessary a list of keywords related to a quick reference of the submitted article.
}
\AbstractPTBR{%
    O resumo deve ser informativo e indicar não só o âmbito geral do artigo, mas também os principais resultados obtidos, os métodos utilizados, o valor do trabalho e as conclusões tiradas. Nenhum número de tabela, número de figura, referência ou expressões matemáticas exibidas devem ser incluídos. Não deve exceder 200 palavras. Após o resumo relacionar as palavras-chave destinadas a uma referência rápida do artigo submetido.
}
\KeywordsENG{keyword 1, keyword 2, keyword 3, keyword 4, keyword 5}
\KeywordsPTBR{palavra-chave 1, palavra-chave 2, palavra-chave  3, palavra-chave  4, palavra-chave  5}


%%%%%%%%%%%%%%%%%%%%%%%%%%%%%%%%%%%%%%%%%%%%%%%
%%    DOCUMENT BEGIN
%%%%%%%%%%%%%%%%%%%%%%%%%%%%%%%%%%%%%%%%%%%%%%%
\begin{document}
\Init

\section*{Introduction}
The Introduction section must identify the nature of the problem under investigation in relation to its antecedents, in a clear and summarized way, while utilizing relevant references. An introduction that consists of an extended literature review is not acceptable.


Types of Articles:

{\begin{enumerate}
	     \item  \textbf{original articles: maximum 15 (fifteen) pages}; 
		 \item  \textbf{review articles: maximum 20 (twenty) pages}; 
\end{enumerate}}

\section*{Materials and methods}
The section Material and methods can precede or follow the section Results and discussion, but they must be separated.

 \section*{Results and discussion}

\begin{enumerate}
   \item Figures and Subtitles:
      \begin{itemize}
         \item Include figures in the text with titles and subtitles, following their appearance order.
         \item Explicitly reference each figure in numerical order (e.g., Figure 1, Figure 2(a), Figure 2(b), etc.).
         \item Place the figure title on top of the figure.
         \item Subtitles should be concise, preferably appearing in a box on the figure or in the title.
      \end{itemize}

   \item File Format for Illustrations:
      \begin{itemize}
         \item Provide a .zip or .tar file containing high-definition (.png format, minimum 300 dpi) illustrations.
         \item Name individual files in the order they appear in the text (e.g., Fig1.png, Fig2.png).
      \end{itemize}

   \item Copyrighted Material:
      \begin{itemize}
         \item If using copyrighted material, obtain written permission from the owners to reproduce the illustrations.
         \item Clearly indicate the source of the illustration and include the phrase "Reproduced from..." in the illustration.
      \end{itemize}

   \item Tables:
      \begin{itemize}
         \item Include tables in the text with explanatory titles for the presented data.
         \item Reference each table in numerical order (e.g., Table 1, Table 2).

         \begin{table}[!htpp]
	% \vspace{0.3cm}
    %\centering
    \caption{Example table}
	% \vspace{-0.6cm
	\label{tab01}
	\renewcommand{\arraystretch}{1.2}
	\begin{center}
    \begin{tabular}{ccc}
    \hline
    \textbf{name } &	\textbf{name  }	& \textbf{name }\\
    \hline		
    A	&23	&- \\
    B	&200	&50\\
    C	&200	&20\\
    D	&400	&10\\
    E	&400	&5\\
    \hline
    \end{tabular}
    \end{center}
    \begin{center}
        \small\textbf{Source:} The authors.
    \end{center}
\end{table}
      \end{itemize}
      

    \vspace{1cm}
   \item Algorithms:
      \begin{itemize}
         \item Present algorithms as text, without enclosing them in boxes.
      \end{itemize}

   \item Math and Equations:
      \begin{itemize}
         \item Prevent confusion between similar characters (e.g., "l" and "1").
         \item Ensure clarity in character placement for superscripts and subscripts.
         \item Use consistent superscript notation for exponential expressions (e.g., 5x10$^3$ instead of 5E03).
         \item Represent vectors and matrices in bold.
         \item The equations are listed sequentially in the text, with right-hand numbering, for example,

\begin{equation}
\label{eq1}
\rho_i=\frac{\sum^{k}_{j=1} d_{i,j}}{(k^2-1)}.
\end{equation}

In equation \eqref{eq1}, ...
      \end{itemize}

      \item The journal \textbf{Semina: Exact and Technological Sciences} from 2023 on utilizes the normative instructions from the American Psychological Association (APA Style - \textcolor{blue}{\href{https://apastyle.apa.org/}{LINK}}) for articles. Check out the Booklet of the Central Library of UEL - \textcolor{blue}{\href{https://sites.uel.br/bibliotecas/wp-content/uploads/2022/03/APOSTILAESTILOAPA.pdf}{LINK}}.
      
      \textbf{Exemples:}

      \parencite{Cirilo2023, Herbert2023}  or  \textcite{Cirilo2023} and  \textcite{Herbert2023}
\end{enumerate}



 
\subsection*{Conclusions}



The Conclusions section, which summarizes the main conclusions of the work, must be presented immediately after the section Results and Discussion.

 




\subsection*{Author contributions}

The journal \textbf{Semina: Exact and Technological Sciences} from 2023 on adopted the system of taxonomy papers for CRediT contributors. \textcolor{blue}{\href{https://ojs.uel.br/revistas/uel/index.php/semexatas/form}{CRediT}}  is a taxonomy of 14 papers that describe the authorial contribution in the making of academic productions. Therefore, all manuscripts submitted to the journal must describe the participation of its authors through a form (\textcolor{blue}{\href{https://ojs.uel.br/revistas/uel/index.php/semexatas/form}{LINK}})


\vspace{-0.05cm}

\subsection*{Conflicts of interest}

\vspace{-0.2cm}
The authors declare no conflict of interest OR  The authors declare......
\vspace{-0.05cm}

\subsection*{Acknowledgments}
 The acknowledgments section in a research paper is optional and allows authors to express gratitude to those who contributed to the research or publication process. It should be concise, clear, and relevant to the work.

 
 


\printbibliography


\end{document}
%%%%%%%%%%%%%%%%%%%%%%%%%%%%%%%%%%%%%%%%%%%%%%%
%%    DOCUMENT END
%%%%%%%%%%%%%%%%%%%%%%%%%%%%%%%%%%%%%%%%%%%%%%%